\section{Discrete Gauss-Bonnet}
\label{sec:discrete}


In this section we explore discrete Gaussian curvature. We give an
application to polygonal mesh processing, state a version of the Guass-Bonnet
theorem of discrete surfaces and provide a proof.

\subsection{Discrete Curvature}

Given a polygonal mesh, each face is planar and therefore has Gaussian curvature zero. Likewise, curvature is not distributed along the edges. Instead, all of the curvature is concentrated at the vertices. At a vertex, however, there is generally no well-defined tangent plane, so the classical definition of Gaussian curvature does not apply directly. Nevertheless, as the triangles of a mesh become arbitrarily small, the mesh increasingly resembles a smooth surface, suggesting that an appropriate notion of discrete curvature should converge to its smooth counterpart.

As in the smooth setting, there are several ways to define discrete Gaussian curvature. One particularly geometric approach uses a discrete analogue of the Gauss map.
Recall that when studying the turning angle of a polygon in the plane (\thmref{simple-bonnet}), we obtained equivalent formulations by measuring either the interior angles or their complementary exterior angles. A similar idea can be used to define curvature at a vertex of a polygonal surface.


Let $v$ be a vertex of a mesh and suppose that $n$ faces are incident to $v$.
Let $\xi_i$ denote the complementary angle associated with the $i^{\text{th}}$ face,
and example is shown in  in \figref{switcheroo}. 

The \emph{discrete Gaussian curvature} at $v$ is defined to be
\begin{definition}[Discrete Curvature]\label{def:discrete-curvature-complement-angle}
$$K(v)=2\pi - \sum_{i=1}^n (\pi - \xi_i)=(2-n)\pi +\sum_{i=1}^n \xi_i.$$
\end{definition}
By \eqnref{sphere-area}, this quantity is equal to the area of a spherical polygon  with interior
angles $\xi_1, \xi_2,\ldots \xi_n.$
If we project the normal vectors on each face onto a sphere and
connect these points with great arcs to create a polygon on the sphere with interior angles
$\xi_1, \xi_2,\ldots \xi_n.$

This interpretation arises for a discrete version of the Gauss map.
For each face incident to $v$, consider its unit normal vector and project that vector
onto the unit sphere. Connecting consecutive normals by great-circle arcs produces
a spherical polygon with interior angles $\xi_1,\xi_2,\ldots,\xi_n.$ The area enclosed by this 
polygon is \defref{discrete-curvature-complement-angle}.


\begin{figure}[htb]
\centering
\includegraphics[width=.25\textwidth]{background/switch-angles}
\caption{The relationship between the angles incident to a vertex and
its complementary angle.}
\label{fig:switcheroo}
\end{figure}

This definition agrees well with our geometric intuition. For example, a point on a flat plane has curvature zero, while a sharper, more pointed vertex has greater curvature than a flatter one. In this sense, the discrete curvature measures how much the surface deviates from being locally flat at a vertex.

Moreover, this definition satisfies a discrete analogue of the Gauss--Bonnet theorem for triangulated surfaces that are homeomorphic to $\Sp^2$. As the spherical polygons associated with all vertices fit together to cover the unit sphere exactly once, the sum of their areas is equal to the area of the sphere, namely $4\pi$. Thus, the total discrete Gaussian curvature of the mesh is 
$$\sum_{v\in V}K(v)=4\pi=2\pi \chi(\Sp^2)$$
which agrees with the Gauss-Bonnet theorem!



\begin{figure}[htb]
\centering
\includegraphics[width=.5\textwidth]{curvature/discrete-curvature}
\caption{Consider the vertex $v$ in the figure on the left. The curvature of $v$
is the area on the sphere shown on the right. We rotated the sphere
in order to see the entire polygon.}
\label{fig:discrete-curvature}
\end{figure}



What about meshes with boundary? 
Consider a triangulated polygonal surface $S$ with boundary $\partial(S)$.
The boundary is a one dimensional piecewise linear curve in $\R^3.$
As with polygons in the plane, at each vertex $v\in \partial(S)$ 
we have an exterior angle.

The interior angle might consist of many triangles.
Let $F_v$  denote the set of faces incident to $v$ and let
$\alpha_f$ denote the angle in face $f$ at $v$.
The \EMPH{discrete geodesic curvature}
of $v$  is
$$k_{g}(v)= \pi-\sum_{f\in F_v}\epsilon_f.$$
See \figref{discrete-geodesic} for an illustration.
Notice that if $v$ lies on a straight line, then $\sum_{i}\epsilon_f=\pi$
and the curvature is zero as we would expect.


\begin{figure}[htb]
\centering
	\includegraphics[width=.25\textwidth]{curvature/discrete-geodesic}
	\caption{The discrete geodesic curvature at a vertex $v$ on the boundary
	of a surface is given by the exterior angle $\epsilon$.}
	\label{fig:discrete-geodesic}
\end{figure}



A version of the Gauss-Bonnet theorem for discrete surfaces is the following.

\begin{theorem}[The Discrete Gauss-Bonnet Theorem] \label{thm:g-b-d}

If $T$ is a triangulated surface with  boundary $\partial T$ then
$$\sum_{v\in V_{int}} K(v) + \sum_{v\in V_{\partial T}} k_g(v) = 2\pi \chi(T)$$
where $K(v)$ is the discrete Gaussian curvature
of a vertex, $k_g(g)$ is the discrete geodesic curvature,  and
$\chi$ is the Euler characteristic.
\end{theorem}

In both the smooth and discrete versions, the Gauss-Bonnet 
theorem tell us that, if we add up curvature
at each vertex the sum will be $2\pi$ times to Euler characteristic.
The curvature at each vertex is computed by considering a `local' neighborhood
around the vertex and we learn global topological information. Thus, the Gauss-Bonnet 
theorem is an example of a local to global principle. 
Conversely, if we know the Euler characteristic we can learn about the curvature
at individual points. 


\subsection{Triangle Cancelation Proof}
\label{sec:proof}


We present a proof of the \thmref{g-b-d}, our proof follows that of Upadhyay \cite{upadhyay2015}.
First, we consider the case where our surface does not have a boundary.
We then extend this case to surfaces with boundary.
\begin{theorem}[Discrete surfaces without boundary]\label{thm:g-b-discete-bdy}
For a triangulated surface $S$ without boundary
$$\sum_{v\in V} K(v)=2\pi \chi(S)$$
where $K(v)$ is the discrete curvature.
\end{theorem}

\begin{proof}

For each vertex $v$ in $S$,
let $deg(v)$ denote the number of edges incident to $v$, let $\alpha_1,\alpha_2,\ldots,\alpha_{\deg{(v)}}$ denote the angles
containing $v$ and let $\xi_i=\pi-\alpha_i$ for each $i$.
By \defref{discrete-curvature-complement-angle}, 
the discrete Gaussian curvature at a $v$ is
 $$K(v)=(2-\deg{(v)})\pi +\sum_{i=1}^{\deg{(v)}} \xi_i.$$
Summing over all vertices in $S$ gives
$$\sum_{v\in V} K(v)=\sum_{v\in V}2\pi - \sum_{v\in V}\deg{(v)}\pi+\sum_{v\in V}\sum_{i=1}^{\deg{(v)}} \xi_i.$$
The first term on the right hand side is $2\pi |V|$. Each edge is incident with two vertices, so the second term is $2\pi |E|$. 
In the third term, we rewrite $\xi_i$ as $\pi-\alpha_i$.

$$ \sum_{v\in V}\sum_{i=1}^{\deg{(v)}} \beta_i= \sum_{v\in V}\sum_{i=1}^{\deg{(v)}} (\pi-\alpha_i).$$
We can reorganize this sum as follows, instead of summing the angles around each vertex we can sum the angles in each face.
Each angle in $S$ is still being counted exactly once. 
Since each face is a triangle, this gives
$$\sum_{v\in V}\sum_{i=1}^{\deg{(v)}} (\pi-\alpha_i)=\sum_{f\in F}\sum_{i=1}^3(\pi-\alpha_i).$$
Since each face is a triangle the sum of the three angles is $\pi$,
so $\sum_{i=1}^3(\pi-\alpha_i)=3\pi-\pi=2\pi.$
Thus, $$\sum_{v\in V} K(v)=2\pi |V|-2\pi |E|+2\pi |F|=2\pi \chi(S)$$ as desired.
\end{proof}

Next, we extend the above proof to the case where $S$ has a boundary
by gluing a copy of $S$ to itself along the boundary.

\begin{theorem}[Discrete surfaces without boundary]\label{thm:g-b-discete}
For a combinatorial surface $S$ with boundary

$$\sum_{v\in S_{\text{int}}} K(v)+\sum_{v\in\partial S}k_g(v)=2\pi \chi(S)$$
where $K(v)$ is the discrete curvature and $k_g(v)$ is the discrete geodesic curvature.
\end{theorem}

\begin{proof}
Take a copy of $S$ and attach it to itself along the boundary.
This creates the surface $2S$ without boundary. Notice,
when we copy $S$ we create two copies of the boundary, and when
we glue we remove one copy of the boundary.
Thus, $$\chi(2S)=2\chi(S)-\chi(\partial S).$$
Since, $\partial S$ is piecewise linear the number of vertices and
edges are equal and there are no faces, so $\chi(\partial S)=0$
and 

\begin{equation} \label{eqn:glue}
\chi(2S)=2\chi(S).
\end{equation}
For $v$ a vertex on the boundary, $k_g(v)$ is half
the discrete Gaussian curvature of $v$ in $2S.$
Thus,

$$\sum_{v\in 2S}K(v)=2\left(\sum_{v\in S_{\text{int}}}K(v)+\sum_{v\in \partial S} k_g(v)\right) =2\pi  \chi(2S).$$
Applying \eqnref{glue},

$$\sum_{v\in S}K(v)+\sum_{v\in \partial S} k_g(v)=2\pi  \chi(S)$$
as desired.

\end{proof}

\subsection{Two Applications}
\label{sec:removing}




We can use a continuous version of the Gauss-Bonnet theorem to derive
the same definition for discrete curvature at a vertex on a polygonal surface.
Note that we need to be carefull not to use the Gauss-Bonnet theorem to
define curvature then use curvature to prove the theorem. 


Suppose we wish to define discrete Gaussian curvature at a vertex
$v$ in a polygonal mesh.
Consider removing the neighborhood around a vertex $v$ consisting
of the faces incident to $v$ and the edges and vertices of these faces call
this $N(v)$.
Using the variable names given in \figref{discrete-geodesic},
the total geodesic curvature

\begin{align}
\int_{\partial N(v)} kg ds &=\sum_i \pi -(\tau_i+\sigma_{i+1})\nonumber  \\
       &=\sum_i \pi -(\tau_i+\sigma_{I}) \nonumber  \\
       &=\sum_i \alpha_i. \nonumber 
\end{align}
Where the second equality comes from the boundary being a circle and reorganizing the sum.
Topologically, this neighborhood is a disk and $\chi(N(v))=1$
so the Gauss-Bonnet theorem tells us that 
the discrete Gaussian curvature at a vertex is
$$K(v)=\int_{N(v)}K da=2\pi-\sum_i \alpha_i.$$
Which is exactly the definition of discrete Gaussian curvature.






As a second application,  definition of discrete Gaussian curvature can be used
to reduce the number of triangles in a mesh. 
Meshes that are obtained by scanning real objects contain noise.
Most meshes that are generated by scanning require a complete
remeshing \cite{remeshing-2003}.
As a first step in remeshing, the curvature at each
vertex needs to be estimated \cite{mmsb-2003}.




Computing the curvature at each vertex in a mesh can allow one to reduce
the number of vertices in the mesh without sacrificing mesh quality.
This leads to simplification process that reduces storage space and improves the efficiency of
 algorithms run on the mesh.
If the curvature at a vertex is zero, then the vertex can be removed,
along with the faces and edges incident to the vertex, without changing 
the mesh.
 If the curvature is small, then few points
are sampled to be included in the re-mesh. If the curvature
is large, then many points are sampled. See \figref{planck-mesh}
for an illustration.

\begin{figure}[htb]
\centering
\includegraphics[width=.7\textwidth]{meshes/planck-mesh.jpeg}
\caption{A mesh of Max Plank. The original mesh is on the left. As we move
from left to right more importance is placed on vertices with large curvature.
Figure from \cite{alliez-2002}.}
\label{fig:planck-mesh}
\end{figure}


The experiments in \cite{mmsb-2003} found that the average
percent error did not exceed $1.3\%$ when using this operator.\todo{what does
this mean?}






\noindent \textbf{Exercises}


\begin{enumerate}
	\item Discrete Gaussian curvature is also called the \emph{angle defect} at a vertex, why?
	
	\item 
	
	\item 
	
\end{enumerate}

\pagebreak